\RequirePackage{fix-cm}
\documentclass[12pt,letterpaper]{article}
\usepackage[margin=.5in]{geometry}
\usepackage{tikz}
\usetikzlibrary{arrows.meta,patterns}
\pdfmapfile{+cm.map}
\pagestyle{empty}
\setlength{\parindent}{0pt}
\begin{document}
\null\begin{tikzpicture}[remember picture,overlay]
\begin{scope}[shift={(current page.north west)},x=1cm,y=-1cm]
\begin{scope}[shift={(1.5,1.5)}]
\node[anchor=north west,inner sep=0,align=left,text width=18.5cm,font=\sffamily\fontsize{11}{13.75}\selectfont] at (0,-0.5) {Week 34 / Hidden turns / Grades 2--3};
\node[anchor=north west,inner sep=0,align=left,text width=17cm,font=\sffamily\fontsize{10}{12.5}\selectfont] at (0,24.7) {Bellingham Math Circle / Week 34 / W34-grades-2-3-v4};
\node[anchor=north west,inner sep=0,align=left,text width=0.5cm,font=\sffamily\fontsize{10}{12.5}\selectfont] at (18,24.7) {1};
\node[anchor=north west,inner sep=0,align=left,text width=18.4cm,font=\sffamily\fontsize{13}{16.25}\selectfont] at (0,0.25) {Fill every spot; neighbors may be the same kind. Only counter kinds and spots count. Copy a pattern on tracing paper with a dot of each counter's color. One partner builds. The other keeps the original still and turns or flips (turns over) the copy to make every dot sit on its own color. A whole turn does not count.};
\draw[gray!45,line width=.5pt] (1.9,6.1000000000000005) circle (1);
\draw[line width=1pt,fill=white] (1.90000,5.10000) circle (0.37);
\fill (1.9,5.1000000000000005) circle (0.1665);
\draw[line width=1pt,fill=white] (2.90000,6.10000) circle (0.37);
\draw[line width=1pt] (2.9,6.1000000000000005) circle (0.1665);
\draw[line width=1pt,fill=white] (1.90000,7.10000) circle (0.37);
\fill (1.9,7.1000000000000005) circle (0.1665);
\draw[line width=1pt,fill=white] (0.90000,6.10000) circle (0.37);
\draw[line width=1pt] (0.8999999999999999,6.1000000000000005) circle (0.1665);
\node[anchor=north,inner sep=0,align=center,text width=4.4cm,font=\sffamily\fontsize{12}{15.0}\selectfont] at (1.9,7.65) {start};
\draw[gray!45,line width=.5pt] (6.7,6.1000000000000005) circle (1);
\draw[line width=1pt,fill=white] (6.70000,5.10000) circle (0.37);
\draw[line width=1pt] (6.7,5.1000000000000005) circle (0.1665);
\draw[line width=1pt,fill=white] (7.70000,6.10000) circle (0.37);
\fill (7.7,6.1000000000000005) circle (0.1665);
\draw[line width=1pt,fill=white] (6.70000,7.10000) circle (0.37);
\draw[line width=1pt] (6.7,7.1000000000000005) circle (0.1665);
\draw[line width=1pt,fill=white] (5.70000,6.10000) circle (0.37);
\fill (5.7,6.1000000000000005) circle (0.1665);
\node[anchor=north,inner sep=0,align=center,text width=4.4cm,font=\sffamily\fontsize{12}{15.0}\selectfont] at (6.7,7.65) {turn 1 spot: no match};
\draw[gray!45,line width=.5pt] (11.5,6.1000000000000005) circle (1);
\draw[line width=1pt,fill=white] (11.50000,5.10000) circle (0.37);
\fill (11.5,5.1000000000000005) circle (0.1665);
\draw[line width=1pt,fill=white] (12.50000,6.10000) circle (0.37);
\draw[line width=1pt] (12.5,6.1000000000000005) circle (0.1665);
\draw[line width=1pt,fill=white] (11.50000,7.10000) circle (0.37);
\fill (11.5,7.1000000000000005) circle (0.1665);
\draw[line width=1pt,fill=white] (10.50000,6.10000) circle (0.37);
\draw[line width=1pt] (10.5,6.1000000000000005) circle (0.1665);
\node[anchor=north,inner sep=0,align=center,text width=4.4cm,font=\sffamily\fontsize{12}{15.0}\selectfont] at (11.5,7.65) {turn 2 spots: match};
\draw[gray!45,line width=.5pt] (16.3,6.1000000000000005) circle (1);
\draw[line width=1pt,fill=white] (16.30000,5.10000) circle (0.37);
\fill (16.3,5.1000000000000005) circle (0.1665);
\draw[line width=1pt,fill=white] (17.30000,6.10000) circle (0.37);
\draw[line width=1pt] (17.3,6.1000000000000005) circle (0.1665);
\draw[line width=1pt,fill=white] (16.30000,7.10000) circle (0.37);
\fill (16.3,7.1000000000000005) circle (0.1665);
\draw[line width=1pt,fill=white] (15.30000,6.10000) circle (0.37);
\draw[line width=1pt] (15.3,6.1000000000000005) circle (0.1665);
\node[anchor=north,inner sep=0,align=center,text width=4.4cm,font=\sffamily\fontsize{12}{15.0}\selectfont] at (16.3,7.65) {flip over the dashed line: match};
\draw[dashed,line width=.8pt] (16.3,4.3500000000000005) -- (16.3,7.8500000000000005);
\draw[-{Stealth},line width=1pt] (3.45,6.1000000000000005) -- (5.15,6.1000000000000005);
\draw[-{Stealth},line width=1pt] (2.62,4.84) .. controls (4.5,3.8) and (8.879999999999999,3.8) .. (10.78,4.84);
\draw[-{Stealth},line width=1pt] (2.62,4.84) .. controls (4.5,3.4000000000000004) and (13.68,3.4000000000000004) .. (15.58,4.84);
\node[anchor=north west,inner sep=0,align=left,text width=18.4cm,font=\sffamily\fontsize{15}{18.75}\selectfont] at (0,9.0) {\textbf{Problem 1:} Find how few kinds of counter can make a pattern on this ring that no turn or flip can match.};
\draw[gray!45,line width=.5pt] (9.2,18.6) circle (5.6);
\draw[line width=1pt,fill=white] (9.20000,13.00000) circle (1.4);
\draw[line width=1pt,fill=white] (14.04974,21.40000) circle (1.4);
\draw[line width=1pt,fill=white] (4.35026,21.40000) circle (1.4);\end{scope}\end{scope}\end{tikzpicture}
\newpage
\null\begin{tikzpicture}[remember picture,overlay]
\begin{scope}[shift={(current page.north west)},x=1cm,y=-1cm]
\begin{scope}[shift={(1.5,1.5)}]
\node[anchor=north west,inner sep=0,align=left,text width=18.5cm,font=\sffamily\fontsize{11}{13.75}\selectfont] at (0,-0.5) {Week 34 / Hidden turns / Grades 2--3};
\node[anchor=north west,inner sep=0,align=left,text width=17cm,font=\sffamily\fontsize{10}{12.5}\selectfont] at (0,24.7) {Bellingham Math Circle / Week 34 / W34-grades-2-3-v4};
\node[anchor=north west,inner sep=0,align=left,text width=0.5cm,font=\sffamily\fontsize{10}{12.5}\selectfont] at (18,24.7) {2};
\node[anchor=north west,inner sep=0,align=left,text width=18.4cm,font=\sffamily\fontsize{15}{18.75}\selectfont] at (0,0.5) {\textbf{Problem 2:} Find how few kinds of counter can make a pattern on this ring that no turn or flip can match.};
\draw[gray!45,line width=.5pt] (9.2,13) circle (7);
\draw[line width=1pt,fill=white] (9.20000,6.00000) circle (1.4);
\draw[line width=1pt,fill=white] (16.20000,13.00000) circle (1.4);
\draw[line width=1pt,fill=white] (9.20000,20.00000) circle (1.4);
\draw[line width=1pt,fill=white] (2.20000,13.00000) circle (1.4);
\draw[gray!40] (0,22.5) -- (18.3,22.5);
\draw[gray!40] (0,24) -- (18.3,24);\end{scope}\end{scope}\end{tikzpicture}
\newpage
\null\begin{tikzpicture}[remember picture,overlay]
\begin{scope}[shift={(current page.north west)},x=1cm,y=-1cm]
\begin{scope}[shift={(1.5,1.5)}]
\node[anchor=north west,inner sep=0,align=left,text width=18.5cm,font=\sffamily\fontsize{11}{13.75}\selectfont] at (0,-0.5) {Week 34 / Hidden turns / Grades 2--3};
\node[anchor=north west,inner sep=0,align=left,text width=17cm,font=\sffamily\fontsize{10}{12.5}\selectfont] at (0,24.7) {Bellingham Math Circle / Week 34 / W34-grades-2-3-v4};
\node[anchor=north west,inner sep=0,align=left,text width=0.5cm,font=\sffamily\fontsize{10}{12.5}\selectfont] at (18,24.7) {3};
\node[anchor=north west,inner sep=0,align=left,text width=18.4cm,font=\sffamily\fontsize{15}{18.75}\selectfont] at (0,0.5) {\textbf{Problem 3:} Find how few kinds of counter can make a pattern on this ring that no turn or flip can match. Can you explain why fewer kinds would fail?};
\draw[gray!45,line width=.5pt] (9.2,13) circle (7);
\draw[line width=1pt,fill=white] (9.20000,6.00000) circle (1.4);
\draw[line width=1pt,fill=white] (15.85740,10.83688) circle (1.4);
\draw[line width=1pt,fill=white] (13.31450,18.66312) circle (1.4);
\draw[line width=1pt,fill=white] (5.08550,18.66312) circle (1.4);
\draw[line width=1pt,fill=white] (2.54260,10.83688) circle (1.4);
\draw[gray!40] (0,23) -- (18.3,23);
\draw[gray!40] (0,24.1) -- (18.3,24.1);\end{scope}\end{scope}\end{tikzpicture}
\newpage
\null\begin{tikzpicture}[remember picture,overlay]
\begin{scope}[shift={(current page.north west)},x=1cm,y=-1cm]
\begin{scope}[shift={(1.5,1.5)}]
\node[anchor=north west,inner sep=0,align=left,text width=18.5cm,font=\sffamily\fontsize{11}{13.75}\selectfont] at (0,-0.5) {Week 34 / Hidden turns / Grades 2--3};
\node[anchor=north west,inner sep=0,align=left,text width=17cm,font=\sffamily\fontsize{10}{12.5}\selectfont] at (0,24.7) {Bellingham Math Circle / Week 34 / W34-grades-2-3-v4};
\node[anchor=north west,inner sep=0,align=left,text width=0.5cm,font=\sffamily\fontsize{10}{12.5}\selectfont] at (18,24.7) {4};
\node[anchor=north west,inner sep=0,align=left,text width=18.4cm,font=\sffamily\fontsize{15}{18.75}\selectfont] at (0,0.5) {\textbf{Problem 4:} Find a six-ring pattern using only two kinds of counter that no turn or flip can match. Explain how you know no turn or flip can match it.};
\draw[gray!45,line width=.5pt] (9.2,13) circle (7);
\draw[line width=1pt,fill=white] (9.20000,6.00000) circle (1.4);
\draw[line width=1pt,fill=white] (15.26218,9.50000) circle (1.4);
\draw[line width=1pt,fill=white] (15.26218,16.50000) circle (1.4);
\draw[line width=1pt,fill=white] (9.20000,20.00000) circle (1.4);
\draw[line width=1pt,fill=white] (3.13782,16.50000) circle (1.4);
\draw[line width=1pt,fill=white] (3.13782,9.50000) circle (1.4);
\draw[gray!40] (0,23) -- (18.3,23);
\draw[gray!40] (0,24.1) -- (18.3,24.1);\end{scope}\end{scope}\end{tikzpicture}
\newpage
\null\begin{tikzpicture}[remember picture,overlay]
\begin{scope}[shift={(current page.north west)},x=1cm,y=-1cm]
\begin{scope}[shift={(1.5,1.5)}]
\node[anchor=north west,inner sep=0,align=left,text width=18.5cm,font=\sffamily\fontsize{11}{13.75}\selectfont] at (0,-0.5) {Week 34 / Hidden turns / Grades 2--3};
\node[anchor=north west,inner sep=0,align=left,text width=17cm,font=\sffamily\fontsize{10}{12.5}\selectfont] at (0,24.7) {Bellingham Math Circle / Week 34 / W34-grades-2-3-v4};
\node[anchor=north west,inner sep=0,align=left,text width=0.5cm,font=\sffamily\fontsize{10}{12.5}\selectfont] at (18,24.7) {5};
\node[anchor=north west,inner sep=0,align=left,text width=18.4cm,font=\sffamily\fontsize{15}{18.75}\selectfont] at (0,0.5) {\textbf{Problem 5:} Find an eight-ring pattern with two kinds of counter that no turn or flip can match. How few counters of one kind can it have?};
\draw[gray!45,line width=.5pt] (9.2,13) circle (7);
\draw[line width=1pt,fill=white] (9.20000,6.00000) circle (1.4);
\draw[line width=1pt,fill=white] (14.14975,8.05025) circle (1.4);
\draw[line width=1pt,fill=white] (16.20000,13.00000) circle (1.4);
\draw[line width=1pt,fill=white] (14.14975,17.94975) circle (1.4);
\draw[line width=1pt,fill=white] (9.20000,20.00000) circle (1.4);
\draw[line width=1pt,fill=white] (4.25025,17.94975) circle (1.4);
\draw[line width=1pt,fill=white] (2.20000,13.00000) circle (1.4);
\draw[line width=1pt,fill=white] (4.25025,8.05025) circle (1.4);
\draw[gray!40] (0,23) -- (18.3,23);
\draw[gray!40] (0,24.1) -- (18.3,24.1);\end{scope}\end{scope}\end{tikzpicture}
\end{document}
